\chapter{Cosmology and the Universe as a Compression}
\index{cosmology}\index{RSVP dynamics}

\section{What Kind of Object Is a Cosmological Model?}

A cosmological model is a compression architecture for physical history. The observable universe contains only a fraction of the information in its past trajectory. The CMB encodes conditions at recombination but not the precise configuration of every field mode. Large-scale structure encodes gravitational amplification of primordial perturbations but not every fluid element's trajectory.

A cosmological theory specifies what information is preserved in the present state, what is discarded, and which historical queries can in principle be answered.

\section{Cosmological Trajectory Spaces}

Four candidates: (1) spacetime-history (complete field histories), (2) cosmological-state (large-scale parameter sequences), (3) accessibility (RSVP field configurations), and (4) observer-conditioned (histories relative to the observer's causal past). The canonical trajectory space is not settled across cosmological theories.

\textbf{Diagnostic Question 1: Affirmative within specific formalisms; contested across them.}

\section{Stability: Diagnostic Question 2}

\textbf{$\Lambda$CDM.} Perturbation theory suggests a summable stability sequence for large-scale observables under temporal interaction order.

\textbf{Inflation.} The trans-Planckian problem \cite{MartinBrandenberger2001, Baumann2009} is a stability question: if Planck-scale physics affects inflationary observables, the stability sequence is not summable under energy-scale interaction order.

\textbf{RSVP.} The Accessibility Relaxation Hypothesis conjectures $\mu(\mathcal{A}_r(F_0)) \leq Ce^{-cr}$ under RSVP dynamics. Three outcomes from Chapter~7 apply at cosmological scale.

\textbf{Diagnostic Question 2: Conditional across all frameworks.}

\section{Cosmological Failure Modes}

\begin{itemize}
  \item \textbf{Trans-Planckian problem} $\leftrightarrow$ stability failure candidate
  \item \textbf{Cosmic variance} $\leftrightarrow$ sufficiency failure (imposed by causal structure)
  \item \textbf{Model underdetermination} $\leftrightarrow$ structural misalignment
  \item \textbf{Information loss behind horizons} $\leftrightarrow$ stability failure or sufficiency failure
  \item \textbf{Dark-sector parameter degeneracy} $\leftrightarrow$ Type II degeneracy
\end{itemize}

\section{RSVP Audit}

RSVP is audited as a candidate theory, not a confirmed one.
\begin{center}\small
\begin{tabular}{lll}
\toprule
Question & Status & Notes \\
\midrule
Decomposition exact? & Yes & Accessibility telescoping \\
Stability summable? & Conditional & Accessibility Relaxation Hypothesis \\
Compression sufficient? & Yes & Skeleton compression \\
Depths aligned? & Conditional & Requires contraction \\
Degeneracy avoided? & Open & Measure-theoretic \\
\bottomrule
\end{tabular}
\end{center}

Open obligations: (1) derive accessibility contraction from RSVP field equations; (2) identify observable signatures distinguishing three outcomes; (3) test inflationary stability against trans-Planckian corrections \cite{Planck2020}; (4) develop compression-theoretic model distinguishability measures; (5) characterize dark-sector degeneracy structure.

\section{The Universe as a Compression of Its Own History}

If cosmological history is compressible (Outcome A), the present state contains recoverable information about nearby alternative histories. If it is not (Outcome C), the limits of cosmological knowledge are structural, not practical. Both possibilities have precise meanings in the framework's language. Both can in principle be distinguished by empirical investigation.
